\documentclass[11pt]{amsart}

\usepackage[T1]{fontenc}
\usepackage{lmodern}
\usepackage{microtype}
\usepackage{mathtools,amssymb,amsthm}
\usepackage[hidelinks]{hyperref}

\hypersetup{
  pdftitle={A 32-digit prime x with x^2+x+1 square-full: a question of Ochem and Zelinsky answered in the affirmative}
}

\newtheorem{theorem}{Theorem}
\theoremstyle{remark}
\newtheorem*{remark}{Remark}
\newtheorem*{question}{Question}

\DeclareMathOperator{\rad}{rad}
\newcommand{\vp}{v}

\title[A 32-digit prime with $x^2+x+1$ square-full]
{A 32-Digit Prime $x$ with $x^2+x+1$ Square-Full:
 A Question of Ochem and Zelinsky Answered in the Affirmative}

\date{}

\subjclass[2020]{11A41, 11A51, 11Y11}
\keywords{square-full number, powerful number, odd perfect number,
primality certificate}

\begin{document}

\begin{abstract}
In the concluding remarks of their paper on odd perfect numbers with exactly one
even exponent greater than $2$, Ochem and Zelinsky ask whether there exists a
prime $x$ such that $x^2+x+1$ is square-full. The answer is yes:
\[
 x=47116128896261596331524418282573
\]
is a prime of $32$ digits with $x^2+x+1=7^3A^2$, where
\[
 A=2544031832644333879196802002161,
\]
and $7\nmid A$, so every prime dividing $x^2+x+1$ divides it at least twice.
The primality of $x$ is certified, not merely tested: $x-1$ is completely
factored, as $2^2\cdot 6299\cdot 9701563\cdot q$ with
$q=192750846182877882539$ prime, and the base $2$ has full order modulo $x$.
The witness is not ours: it is $a(27)$ of the OEIS entry A296376, whose
underlying family that entry traces to a note of Majol from 1902.
\end{abstract}

\maketitle

\section{The question}

An integer $n\ge 1$ is \emph{square-full} (equivalently \emph{powerful}) if
$(\rad n)^2\mid n$, that is, if every prime dividing $n$ divides it at least
twice; this is the definition used in \cite{OZ}. In the concluding remarks of
\cite{OZ}, Ochem and Zelinsky pose:

\begin{question}[\cite{OZ}, Question~9]
Does there exist a prime $x$ such that $x^2+x+1$ is square-full?
\end{question}

\noindent
They pose it because a negative answer would reduce the exponent pairs $(a,b)$
admissible in their analysis; the affirmative answer below rules out obtaining
that reduction from this question. Note that ``square-full'' is neither
``squarefree'' nor ``a perfect power''.

\section{The answer}

\begin{theorem}\label{thm:main}
Let
\[
 x=47116128896261596331524418282573 .
\]
Then $x$ is prime and $x^2+x+1$ is square-full. Consequently the question above
has the answer \emph{yes}.
\end{theorem}

Everything a reader needs is the integer $x$ above together with the following
three decimal strings. Write $n=x^2+x+1$ and
{\small
\begin{equation}\label{eq:n}
 n=2219929602169136991765718912586970986546248112210754339293782903,
\end{equation}}
\begin{align}
 A&=2544031832644333879196802002161, \label{eq:A}\\
 x-1&=2^2\cdot 6299\cdot 9701563\cdot 192750846182877882539. \label{eq:xm1}
\end{align}

\begin{proof}[Proof that $n$ is square-full]
Two multiplications suffice. First,
\begin{equation}\label{eq:key}
 n=7^3A^2,
\end{equation}
which is \eqref{eq:n} against \eqref{eq:A}; and second $A\equiv 3\pmod 7$, so
$7\nmid A$. Now let $p$ be any prime dividing $n$. If $p=7$ then
$\vp_7(n)=3+2\vp_7(A)=3\ge2$. If $p\neq7$ then $p\mid A^2$ by \eqref{eq:key},
hence $p\mid A$ and $\vp_p(n)=2\vp_p(A)\ge2$. Every exponent in $n$ is therefore
at least $2$, i.e.\ $(\rad n)^2\mid n$.
\end{proof}

No factorisation of $A$ is needed for that argument, which is why it is
checkable by hand. The complete factorisation, which the argument does not use,
is
\begin{equation}\label{eq:full}
\begin{aligned}
 n&=7^3\cdot 43^2\cdot 3416683^2\cdot 12353161^2\cdot 1401752363863729^2,\\
 A&=43\cdot 3416683\cdot 12353161\cdot 1401752363863729 ,
\end{aligned}
\end{equation}
each listed factor being prime: $43$, $3416683$ and $12353161$ by trial
division, and $1401752363863729$ by the certificate of the next proof applied
with base $7$ to
\[
 1401752363863729-1=2^4\cdot3\cdot37\cdot107\cdot7376401679
\]
($7376401679$ prime by trial division; base $2$ fails there, being a quadratic
residue modulo $1401752363863729$).

\begin{proof}[Proof that $x$ is prime]
This is the converse of Fermat's theorem in the form of Lucas, as in
\cite[Theorem~1]{BLS}: if $a^{x-1}\equiv1\pmod x$ and
$a^{(x-1)/q}\not\equiv1\pmod x$ for every prime $q\mid x-1$, then $x$ is prime.
Take $a=2$. The factorisation \eqref{eq:xm1} of $x-1$ is complete, and
$2^{x-1}\equiv1\pmod x$, while the four quotient powers are
{\small
\begin{align*}
 2^{(x-1)/2}          &\equiv 47116128896261596331524418282572 \pmod x,\\
 2^{(x-1)/6299}       &\equiv 42436530233755792060895430447505 \pmod x,\\
 2^{(x-1)/9701563}    &\equiv 21277589209484664623743014586336 \pmod x,\\
 2^{(x-1)/192750846182877882539}
                      &\equiv 19695209395513058587964373225481 \pmod x,
\end{align*}}
none of them $1$. It remains to know that the four factors in \eqref{eq:xm1} are
prime. For $2$, $6299$ and $9701563$ this is trial division. For the largest
factor $m$ we give the same certificate one level down: with
\[
 m=192750846182877882539,\qquad m-1=2\cdot 17^2\cdot 27779\cdot 12004714807399,
\]
again completely factored, the base $2$ satisfies $2^{m-1}\equiv1\pmod m$ and
$2^{(m-1)/q}\not\equiv1\pmod m$ for each of the four primes
$q\mid m-1$; and $17$, $27779$ and $12004714807399$ are
prime by trial division, the largest of them needing divisors only up to
$3464782$. The certificate is therefore grounded in exhaustive trial division at
every leaf, and no probable-prime test enters.
\end{proof}

Theorem~\ref{thm:main} follows.

\section{Provenance of the witness}

The integer is not new. It is $a(27)$ of OEIS A296376, ``Natural numbers $x$
such that $7y^2=x^2+x+1$ has a solution in natural numbers'' \cite{A296376},
contributed by Jeffrey Shallit in December 2017, and it occurs as line~27 of the
published $b$-file \texttt{b296376.txt} of that entry, computed by Colin Barker;
the entry traces the family to a note of Majol from 1902 \cite{Majol}. Which
members of the family are square-full is immediate from the defining property
rather than new: if $x=a(k)$ then $x^2+x+1=7y^2$, and such a value is
square-full precisely when $7\mid y$, since $\vp_7(x^2+x+1)=1+2\vp_7(y)$ while
every other exponent of $7y^2$ is even. No priority is claimed for the
observation that this particular member is prime; our literature search was not
exhaustive.

\section{Verification}

The file \texttt{verify.py} accompanying this note, with its recorded run in
\texttt{verify.output.txt}, takes the decimal strings
\eqref{eq:n}--\eqref{eq:xm1} and re-derives the mathematical assertions of
Section~2: the identity $n=7^3A^2$ with $7\nmid A$; the factorisation
\eqref{eq:full}, with each listed prime proved prime, each exponent checked to
be at least $2$, and $(\rad n)^2\mid n$ verified directly; and the Lucas
certificate for $x$ together with its subsidiary certificates, every leaf
grounded in exhaustive trial division. It also runs controls in both directions
on its square-full decider and a census deciding every $x$ with $1\le x\le10^8$,
where the only values with $x^2+x+1$ square-full are the composite $x=18$ and
$x=88916$. No minimality is claimed for the witness: nothing here excludes a
prime witness with between $9$ and $31$ digits. Nothing bibliographic is
re-derived --- the statements of Section~3 are transcribed, as the program's
closing scope note, reproduced in the accompanying \texttt{REVIEW\_NOTE.md},
records. The program uses exact integer arithmetic throughout, no floating-point
comparison, no randomness, and the Python standard library only, and reports
\begin{center}
\texttt{VERDICT: ALL 74 CHECKS PASS}
\end{center}
Every input it consumes is printed above. The decisive steps are elementary
without being a single line of arithmetic: \eqref{eq:key} is one multiplication
and the certificate for $x$ is five modular exponentiations, to which grounding
that certificate adds the subsidiary certificate for $192750846182877882539$
and the trial divisions at its leaves.

\begin{thebibliography}{9}

\bibitem{BLS}
J.~Brillhart, D.~H. Lehmer, and J.~L. Selfridge,
\emph{New primality criteria and factorizations of $2^m\pm1$},
Math. Comp. \textbf{29} (1975), 620--647.
\href{https://doi.org/10.1090/S0025-5718-1975-0384673-1}%
{doi:10.1090/S0025-5718-1975-0384673-1}.

\bibitem{Majol}
E.-A. Majol,
Note \#2228,
\emph{L'Interm\'ediaire des Math\'ematiciens} \textbf{9} (1902), 183--185.
Cited as the reference of \cite{A296376}.

\bibitem{OZ}
P.~Ochem and J.~Zelinsky,
\emph{On odd perfect numbers with exactly one even exponent greater than $2$},
arXiv:2607.19746v1 [math.GM], 22 July 2026.
\href{https://arxiv.org/abs/2607.19746}{arXiv:2607.19746}.
The question reproduced above is Question~9, in the concluding remarks of that
e-print; the numbering of a printed version, if any, may differ, which is why
the question is quoted in full.

\bibitem{A296376}
OEIS Foundation Inc., sequence A296376,
\emph{Natural numbers $x$ such that $7y^2=x^2+x+1$ has a solution in natural
numbers}, contributed by J.~Shallit, 11 December 2017; $b$-file
\texttt{b296376.txt} ($800$ terms) by C.~Barker.
\href{https://oeis.org/A296376}{oeis.org/A296376}.

\end{thebibliography}

\end{document}
